\documentclass[
  man,
  longtable,
  nolmodern,
  notxfonts,
  notimes,
  colorlinks=true,linkcolor=blue,citecolor=blue,urlcolor=blue]{apa7}

\usepackage{amsmath}
\usepackage{amssymb}




\RequirePackage{longtable}
\RequirePackage{threeparttablex}

\makeatletter
\renewcommand{\paragraph}{\@startsection{paragraph}{4}{\parindent}%
	{0\baselineskip \@plus 0.2ex \@minus 0.2ex}%
	{-.5em}%
	{\normalfont\normalsize\bfseries\typesectitle}}

\renewcommand{\subparagraph}[1]{\@startsection{subparagraph}{5}{0.5em}%
	{0\baselineskip \@plus 0.2ex \@minus 0.2ex}%
	{-\z@\relax}%
	{\normalfont\normalsize\bfseries\itshape\hspace{\parindent}{#1}\textit{\addperi}}{\relax}}
\makeatother




\usepackage{longtable, booktabs, multirow, multicol, colortbl, hhline, caption, array, float, xpatch}
\usepackage{subcaption}


\renewcommand\thesubfigure{\Alph{subfigure}}
\setcounter{topnumber}{2}
\setcounter{bottomnumber}{2}
\setcounter{totalnumber}{4}
\renewcommand{\topfraction}{0.85}
\renewcommand{\bottomfraction}{0.85}
\renewcommand{\textfraction}{0.15}
\renewcommand{\floatpagefraction}{0.7}

\usepackage{tcolorbox}
\tcbuselibrary{listings,theorems, breakable, skins}
\usepackage{fontawesome5}

\definecolor{quarto-callout-color}{HTML}{909090}
\definecolor{quarto-callout-note-color}{HTML}{0758E5}
\definecolor{quarto-callout-important-color}{HTML}{CC1914}
\definecolor{quarto-callout-warning-color}{HTML}{EB9113}
\definecolor{quarto-callout-tip-color}{HTML}{00A047}
\definecolor{quarto-callout-caution-color}{HTML}{FC5300}
\definecolor{quarto-callout-color-frame}{HTML}{ACACAC}
\definecolor{quarto-callout-note-color-frame}{HTML}{4582EC}
\definecolor{quarto-callout-important-color-frame}{HTML}{D9534F}
\definecolor{quarto-callout-warning-color-frame}{HTML}{F0AD4E}
\definecolor{quarto-callout-tip-color-frame}{HTML}{02B875}
\definecolor{quarto-callout-caution-color-frame}{HTML}{FD7E14}

%\newlength\Oldarrayrulewidth
%\newlength\Oldtabcolsep


\usepackage{hyperref}



\usepackage{color}
\usepackage{fancyvrb}
\newcommand{\VerbBar}{|}
\newcommand{\VERB}{\Verb[commandchars=\\\{\}]}
\DefineVerbatimEnvironment{Highlighting}{Verbatim}{commandchars=\\\{\}}
% Add ',fontsize=\small' for more characters per line
\usepackage{framed}
\definecolor{shadecolor}{RGB}{241,243,245}
\newenvironment{Shaded}{\begin{snugshade}}{\end{snugshade}}
\newcommand{\AlertTok}[1]{\textcolor[rgb]{0.68,0.00,0.00}{#1}}
\newcommand{\AnnotationTok}[1]{\textcolor[rgb]{0.37,0.37,0.37}{#1}}
\newcommand{\AttributeTok}[1]{\textcolor[rgb]{0.40,0.45,0.13}{#1}}
\newcommand{\BaseNTok}[1]{\textcolor[rgb]{0.68,0.00,0.00}{#1}}
\newcommand{\BuiltInTok}[1]{\textcolor[rgb]{0.00,0.23,0.31}{#1}}
\newcommand{\CharTok}[1]{\textcolor[rgb]{0.13,0.47,0.30}{#1}}
\newcommand{\CommentTok}[1]{\textcolor[rgb]{0.37,0.37,0.37}{#1}}
\newcommand{\CommentVarTok}[1]{\textcolor[rgb]{0.37,0.37,0.37}{\textit{#1}}}
\newcommand{\ConstantTok}[1]{\textcolor[rgb]{0.56,0.35,0.01}{#1}}
\newcommand{\ControlFlowTok}[1]{\textcolor[rgb]{0.00,0.23,0.31}{\textbf{#1}}}
\newcommand{\DataTypeTok}[1]{\textcolor[rgb]{0.68,0.00,0.00}{#1}}
\newcommand{\DecValTok}[1]{\textcolor[rgb]{0.68,0.00,0.00}{#1}}
\newcommand{\DocumentationTok}[1]{\textcolor[rgb]{0.37,0.37,0.37}{\textit{#1}}}
\newcommand{\ErrorTok}[1]{\textcolor[rgb]{0.68,0.00,0.00}{#1}}
\newcommand{\ExtensionTok}[1]{\textcolor[rgb]{0.00,0.23,0.31}{#1}}
\newcommand{\FloatTok}[1]{\textcolor[rgb]{0.68,0.00,0.00}{#1}}
\newcommand{\FunctionTok}[1]{\textcolor[rgb]{0.28,0.35,0.67}{#1}}
\newcommand{\ImportTok}[1]{\textcolor[rgb]{0.00,0.46,0.62}{#1}}
\newcommand{\InformationTok}[1]{\textcolor[rgb]{0.37,0.37,0.37}{#1}}
\newcommand{\KeywordTok}[1]{\textcolor[rgb]{0.00,0.23,0.31}{\textbf{#1}}}
\newcommand{\NormalTok}[1]{\textcolor[rgb]{0.00,0.23,0.31}{#1}}
\newcommand{\OperatorTok}[1]{\textcolor[rgb]{0.37,0.37,0.37}{#1}}
\newcommand{\OtherTok}[1]{\textcolor[rgb]{0.00,0.23,0.31}{#1}}
\newcommand{\PreprocessorTok}[1]{\textcolor[rgb]{0.68,0.00,0.00}{#1}}
\newcommand{\RegionMarkerTok}[1]{\textcolor[rgb]{0.00,0.23,0.31}{#1}}
\newcommand{\SpecialCharTok}[1]{\textcolor[rgb]{0.37,0.37,0.37}{#1}}
\newcommand{\SpecialStringTok}[1]{\textcolor[rgb]{0.13,0.47,0.30}{#1}}
\newcommand{\StringTok}[1]{\textcolor[rgb]{0.13,0.47,0.30}{#1}}
\newcommand{\VariableTok}[1]{\textcolor[rgb]{0.07,0.07,0.07}{#1}}
\newcommand{\VerbatimStringTok}[1]{\textcolor[rgb]{0.13,0.47,0.30}{#1}}
\newcommand{\WarningTok}[1]{\textcolor[rgb]{0.37,0.37,0.37}{\textit{#1}}}

\providecommand{\tightlist}{%
  \setlength{\itemsep}{0pt}\setlength{\parskip}{0pt}}
\usepackage{longtable,booktabs,array}
\usepackage{calc} % for calculating minipage widths
% Correct order of tables after \paragraph or \subparagraph
\usepackage{etoolbox}
\makeatletter
\patchcmd\longtable{\par}{\if@noskipsec\mbox{}\fi\par}{}{}
\makeatother
% Allow footnotes in longtable head/foot
\IfFileExists{footnotehyper.sty}{\usepackage{footnotehyper}}{\usepackage{footnote}}
\makesavenoteenv{longtable}

\usepackage{graphicx}
\makeatletter
\newsavebox\pandoc@box
\newcommand*\pandocbounded[1]{% scales image to fit in text height/width
  \sbox\pandoc@box{#1}%
  \Gscale@div\@tempa{\textheight}{\dimexpr\ht\pandoc@box+\dp\pandoc@box\relax}%
  \Gscale@div\@tempb{\linewidth}{\wd\pandoc@box}%
  \ifdim\@tempb\p@<\@tempa\p@\let\@tempa\@tempb\fi% select the smaller of both
  \ifdim\@tempa\p@<\p@\scalebox{\@tempa}{\usebox\pandoc@box}%
  \else\usebox{\pandoc@box}%
  \fi%
}
% Set default figure placement to htbp
\def\fps@figure{htbp}
\makeatother


% definitions for citeproc citations
\NewDocumentCommand\citeproctext{}{}
\NewDocumentCommand\citeproc{mm}{%
  \begingroup\def\citeproctext{#2}\cite{#1}\endgroup}
\makeatletter
 % allow citations to break across lines
 \let\@cite@ofmt\@firstofone
 % avoid brackets around text for \cite:
 \def\@biblabel#1{}
 \def\@cite#1#2{{#1\if@tempswa , #2\fi}}
\makeatother
\newlength{\cslhangindent}
\setlength{\cslhangindent}{1.5em}
\newlength{\csllabelwidth}
\setlength{\csllabelwidth}{3em}
\newenvironment{CSLReferences}[2] % #1 hanging-indent, #2 entry-spacing
 {\begin{list}{}{%
  \setlength{\itemindent}{0pt}
  \setlength{\leftmargin}{0pt}
  \setlength{\parsep}{0pt}
  % turn on hanging indent if param 1 is 1
  \ifodd #1
   \setlength{\leftmargin}{\cslhangindent}
   \setlength{\itemindent}{-1\cslhangindent}
  \fi
  % set entry spacing
  \setlength{\itemsep}{#2\baselineskip}}}
 {\end{list}}
\usepackage{calc}
\newcommand{\CSLBlock}[1]{\hfill\break\parbox[t]{\linewidth}{\strut\ignorespaces#1\strut}}
\newcommand{\CSLLeftMargin}[1]{\parbox[t]{\csllabelwidth}{\strut#1\strut}}
\newcommand{\CSLRightInline}[1]{\parbox[t]{\linewidth - \csllabelwidth}{\strut#1\strut}}
\newcommand{\CSLIndent}[1]{\hspace{\cslhangindent}#1}





\usepackage{newtx}

\defaultfontfeatures{Scale=MatchLowercase}
\defaultfontfeatures[\rmfamily]{Ligatures=TeX,Scale=1}





\title{Fitting the Memory Measurement Model (M3) with the bmm R Package:
A Tutorial}


\shorttitle{Fitting the M3 with the bmm R Package}


\usepackage{etoolbox}









\authorsnames[{1},{1},{1,2}]{Gidon T. Frischkorn,Isabel Courage,Chenyu
Li}







\authorsaffiliations{
{Department of Psychology, University of Zurich},{Department of
Psychology, University of Chicago}}




\leftheader{Frischkorn, Courage and Li}



\abstract{The memory measurement model (M3) decomposes categorical
response data from memory tasks into latent activation parameters,
separating item memory from context-specific binding via competitive
retrieval. Fitting M3 has so far required custom JAGS or Stan code for
each task design. The bmm R package provides a standardized interface to
specify, fit, and evaluate Bayesian hierarchical M3 with minimal code.
This tutorial demonstrates the bmm workflow through three worked
examples: fitting a simple span M3 to cued recall data, where Bayes
factor model comparison favors the softmax over the simple choice rule;
extending the workflow to a complex span task with distractor categories
and a filtering parameter; and defining a fully custom M3 for a memory
updating task, validating it through parameter recovery across a range
of sample sizes and trial counts. Complete R scripts accompany each
tutorial. After completing the tutorials, readers will be able to fit
standard and custom M3 to their own data and validate new model
specifications through simulation. }

\keywords{working memory, measurement model, multinomial model, Bayesian
estimation, R package}

\authornote{\par{\addORCIDlink{Gidon T.
Frischkorn}{0000-0002-5055-9764}}\par{\addORCIDlink{Isabel
Courage}{0009-0000-8493-2240}}\par{\addORCIDlink{Chenyu
Li}{0000-0001-8815-6189}} 

\par{This study was not preregistered. Data, analysis code, and
companion R scripts for all three tutorials are available at
{[}repository link{]}.  The authors declare no conflicts of
interest. {[}FUNDING INFORMATION{]} {[}ACKNOWLEDGMENTS{]} Gidon T.
Frischkorn: Conceptualization, Methodology, Software, Formal Analysis,
Visualization, Writing -- Original Draft. Isabel Courage: Software,
Formal Analysis, Writing -- Review \& Editing. Chenyu Li: Software, Data
Curation, Writing -- Review \& Editing. }
\par{Correspondence concerning this article should be addressed to Gidon
T.
Frischkorn, Email: \href{mailto:gidon.frischkorn@psychologie.uzh.ch}{gidon.frischkorn@psychologie.uzh.ch}}
}

\makeatletter
\let\endoldlt\endlongtable
\def\endlongtable{
\hline
\endoldlt
}
\makeatother
\RequirePackage{longtable}
\DeclareDelayedFloatFlavor{longtable}{table}

\urlstyle{same}



\makeatletter
\@ifpackageloaded{caption}{}{\usepackage{caption}}
\AtBeginDocument{%
\ifdefined\contentsname
  \renewcommand*\contentsname{Table of contents}
\else
  \newcommand\contentsname{Table of contents}
\fi
\ifdefined\listfigurename
  \renewcommand*\listfigurename{List of Figures}
\else
  \newcommand\listfigurename{List of Figures}
\fi
\ifdefined\listtablename
  \renewcommand*\listtablename{List of Tables}
\else
  \newcommand\listtablename{List of Tables}
\fi
\ifdefined\figurename
  \renewcommand*\figurename{Figure}
\else
  \newcommand\figurename{Figure}
\fi
\ifdefined\tablename
  \renewcommand*\tablename{Table}
\else
  \newcommand\tablename{Table}
\fi
}
\@ifpackageloaded{float}{}{\usepackage{float}}
\floatstyle{ruled}
\@ifundefined{c@chapter}{\newfloat{codelisting}{h}{lop}}{\newfloat{codelisting}{h}{lop}[chapter]}
\floatname{codelisting}{Listing}
\newcommand*\listoflistings{\listof{codelisting}{List of Listings}}
\makeatother
\makeatletter
\makeatother
\makeatletter
\@ifpackageloaded{caption}{}{\usepackage{caption}}
\@ifpackageloaded{subcaption}{}{\usepackage{subcaption}}
\makeatother

% From https://tex.stackexchange.com/a/645996/211326
%%% apa7 doesn't want to add appendix section titles in the toc
%%% let's make it do it
\makeatletter
\xpatchcmd{\appendix}
  {\par}
  {\addcontentsline{toc}{section}{\@currentlabelname}\par}
  {}{}
\makeatother

%% Disable longtable counter
%% https://tex.stackexchange.com/a/248395/211326

\usepackage{etoolbox}

\makeatletter
\patchcmd{\LT@caption}
  {\bgroup}
  {\bgroup\global\LTpatch@captiontrue}
  {}{}
\patchcmd{\longtable}
  {\par}
  {\par\global\LTpatch@captionfalse}
  {}{}
\apptocmd{\endlongtable}
  {\ifLTpatch@caption\else\addtocounter{table}{-1}\fi}
  {}{}
\newif\ifLTpatch@caption
\makeatother

\begin{document}

\maketitle




\setlength\LTleft{0pt}




\section{Introduction}\label{introduction}

Observed performance in memory and decision tasks is often summarized by
overall accuracy. Yet, accuracy collapses across qualitatively distinct
response categories that may arise from different latent processes.
Specific error types are diagnostically informative because different
cognitive processes imply different conditional response probabilities
over observable outcomes. Multinomial processing tree (MPT) models
formalize this logic by treating categorical responses---including
distinct error classes---as consequences of latent states such as
successful retrieval, guessing, or source confusion
(\citeproc{ref-batchelder1999theoretical}{Batchelder \& Riefer, 1999};
\citeproc{ref-erdfelder2009multinomial}{Erdfelder et al., 2009}).
Similarly, process-dissociation approaches demonstrate that patterns of
correct and incorrect responses across task constraints can separate the
contributions of multiple processes within a single task
(\citeproc{ref-jacoby1991process}{Jacoby, 1991};
\citeproc{ref-yonelinas2012process}{Yonelinas \& Jacoby, 2012}). Across
these perspectives, the full multinomial response distribution -- not
accuracy alone -- carries more information about latent cognitive
mechanisms than a single accuracy score.

Working memory research frequently relies on tasks that inherently yield
response frequencies across discrete outcome categories, such as correct
responses, feature confusions, location errors, or intrusions
(\citeproc{ref-oberauer2018benchmarks}{Oberauer et al., 2018}). These
data are naturally multinomial: each trial produces one of several
mutually exclusive response types. Theoretical claims can often be
evaluated by how experimental manipulations and individual differences
redistribute probability mass across these categories. A multinomial
measurement perspective therefore treats the response profile as the
observable consequence of latent memory states and retrieval processes.

The recently proposed memory measurement model {[}M3; Oberauer and
Lewandowsky (\citeproc{ref-oberauer2019simple}{2019}){]} framework
provides one such decomposition. Its core architecture assumes that
multiple activation sources contribute to each response category and
that selection among categories follows a competitive retrieval rule
(\citeproc{ref-luce1959individual}{Luce, 1959}). Although M3 was
introduced for working memory tasks (decomposing performance into item
activation and context-specific binding strength), this architecture is
domain-general: whenever a task yields categorical responses that can be
traced to distinct latent activation sources driving competitive
selection, the framework applies.

However, because activation formulas vary across tasks and response
categories, each new application requires translating its specific model
into custom statistical code. Accordingly, fitting M3 so far required
custom JAGS or Stan code, which made each implementation error-prone and
hard to reuse across designs. The \texttt{bmm} R package
(\citeproc{ref-bmm2025}{Popov et al., 2025}) addresses this limitation
by providing a standardized workflow to specify, fit, and evaluate
Bayesian hierarchical M3 with minimal code, while remaining flexible
enough to accommodate new task designs and custom model specifications.
In this tutorial, we illustrate this approach with three worked examples
covering the standard simple span (ss) and complex span (cs) tasks, and
a custom extension.

\section{The M3 Framework}\label{the-m3-framework}

The M3 framework is designed for categorical response data, tasks in
which each trial produces one of \(K\) mutually exclusive response
types. Rather than collapsing responses into a single accuracy score
(correct vs.~all other categories), M3 treats the full frequency
distribution across categories as the outcome to be explained. The model
decomposes observed response proportions into latent activation
parameters that reflect distinct cognitive processes contributing to
memory performance.

Consider a cued recall task in which a participant sees a probe and
selects their response from a set containing the correct target, other
studied items, and new lures. From a memory perspective, each response
falls into one of three categories: \emph{correct}, \emph{other-list
intrusion}, or \emph{not-presented lure} (NPL). By specifying activation
formulas for each category, M3 explains why each category is selected
with a certain frequency: different activation sources feed into each
category, and response probabilities follow from competitive selection.

\subsection{Activation Sources}\label{activation-sources}

In the M3 framework, each response category receives activation from one
or more latent sources. For the cued recall task introduced above, these
parameters are:

\begin{itemize}
\tightlist
\item
  \(b\): \textbf{baseline activation}, the tendency to select any
  response option in the absence of memory-based evidence. This
  parameter is fixed (not estimated; e.g., \(b = 0\) for the softmax
  rule) to set the scale of the model.
\item
  \(a\): \textbf{item activation}, or general memory strength for
  studied items. An item that was encountered during study receives this
  activation regardless of its position or context.
\item
  \(c\): \textbf{context binding}, the strength of the association
  between an item and its specific study context (e.g., serial position,
  paired associate).
\end{itemize}

This so-called simple span model
(\citeproc{ref-oberauer2019simple}{Oberauer \& Lewandowsky, 2019})
assumes that these activation sources combine additively to produce
total activations for each response category:

\[
\begin{aligned}
\text{act}_{\text{correct}} &= b + a + c \\
\text{act}_{\text{other}} &= b + a \\
\text{act}_{\text{NPL}} &= b
\end{aligned}
\]

The correct target receives activation from all three sources. Other
studied items receive item activation but lack the context binding to
the probed position. Lures receive only baseline activation.

\subsection{Competitive Selection and Choice
Rules}\label{competitive-selection-and-choice-rules}

The activation formulas specify the activation of a single item in each
category. The number of response options per category (\(n_k\)) also
matters. If the recognition set contains 1 correct target, 3 other
studied items, and 4 lures, the model weights each category's activation
by its number of available options before computing response
probabilities.

M3 provides two rules for converting activations into response
probabilities. Both implement competitive selection, where a category's
probability depends on its activation relative to the total, but they
differ in the transformation applied.

The \emph{simple} (linear) choice rule
(\citeproc{ref-luce1959individual}{Luce, 1959}) normalizes activations
directly:

\[P(k) = \frac{n_k \cdot \text{act}_k}{\sum_{j=1}^{K} n_j \cdot \text{act}_j}\]

where \(\text{act}_k\) is the total activation for category \(k\) and
\(n_k\) is the number of response options in that category. This rule
requires all activations to be positive, so the baseline \(b\) is fixed
at a small positive value (0.1 by default).

The \emph{softmax} (exponential) choice rule applies an exponential
transformation before normalization:

\[P(k) = \frac{n_k \cdot \exp(\text{act}_k)}{\sum_{j=1}^{K} n_j \cdot \exp(\text{act}_j)}\]

Because \(\exp(x) > 0\) for all \(x\), softmax guarantees valid
probabilities regardless of activation values. The baseline is fixed at
\(b = 0\) (since \(\exp(0) = 1\), lures receive a positive baseline
contribution automatically).

The softmax rule is formally equivalent to an \(n\)-alternative
forced-choice signal detection model with Gumbel-distributed noise
(\citeproc{ref-oberauer2019simple}{Oberauer \& Lewandowsky, 2019};
\citeproc{ref-thurstone1927law}{Thurstone, 1927};
\citeproc{ref-yellott1977relationship}{Yellott, 1977}). Under this
interpretation, each category's activation represents a deterministic
signal, independent Gumbel noise is added to each alternative, and the
observer selects the category with the highest noisy value. This
equivalence gives activation parameters a principled interpretation as
discriminability measures. Recent empirical comparisons have favored the
softmax rule over the simple rule for working memory data
(\citeproc{ref-li2026filtering}{Li et al., in press};
\citeproc{ref-oberauer2023measurement}{Oberauer, 2023}). Tutorial 1
compares both rules; Tutorials 2 and 3 use the softmax rule as the
default.

\subsection{Model Versions}\label{model-versions}

The \texttt{bmm} package provides three M3 versions that differ in the
number of response categories and estimated parameters.

\begin{itemize}
\item
  \textbf{Simple span} (\texttt{ss}): Three categories (correct, other,
  NPL) with two free parameters (\(a\), \(c\)). Suitable for tasks
  without distractor processing, such as cued recall from a single study
  list.
\item
  \textbf{Complex span} (\texttt{cs}): Five categories --- correct,
  paired distractor, other memorandum, other distractor, and NPL ---
  with three free parameters (\(a\), \(c\), \(f\)). The additional
  parameter \(f\) captures distractor filtering efficiency on a scale
  from 0 to 1: \(f = 0\) indicates perfect filtering (distractors
  receive only baseline activation), while \(f = 1\) indicates no
  filtering (distractors are activated as strongly as memoranda).
\item
  \textbf{Custom}: User-defined response categories, activation
  formulas, and parameters. This version allows researchers to specify
  M3 of arbitrary complexity for novel task designs, as long as the
  model remains identified.
\end{itemize}

\subsection{The bmm Implementation}\label{the-bmm-implementation}

Fitting M3 with the \texttt{bmm} R package (\citeproc{ref-bmm2025}{Popov
et al., 2025}) follows a three-step workflow:

\begin{enumerate}
\def\labelenumi{\arabic{enumi}.}
\tightlist
\item
  \textbf{Define the model} with \texttt{m3()}: specify response
  categories, the number of response options per category, the model
  version, and the choice rule.
\item
  \textbf{Specify predictor formulas} with \texttt{bmf()}: define how
  experimental conditions and individual differences predict each
  activation parameter, using the same formula syntax as \texttt{brms}
  (\citeproc{ref-burkner2017brms}{Bürkner, 2017}).
\item
  \textbf{Fit the model} with \texttt{bmm()}: estimate parameters via
  Bayesian inference using Stan
  (\citeproc{ref-carpenter2017stan}{Carpenter et al., 2017}), with full
  support for hierarchical (multilevel) models, prior specification, and
  posterior diagnostics.
\end{enumerate}

This workflow replaces the previous approach of writing custom JAGS code
for each new experiment, in which the researcher had to manually specify
the likelihood, priors, and sampling scheme --- a process that was
error-prone and difficult to adapt across designs. Readers familiar with
multinomial processing tree (MPT) models will note similarities: both
frameworks decompose categorical response distributions into latent
process parameters. However, whereas MPT models use discrete processing
trees with binary branching probabilities
(\citeproc{ref-batchelder1999theoretical}{Batchelder \& Riefer, 1999};
e.g., \citeproc{ref-heck2018treebugs}{Heck et al., 2018} for Bayesian
hierarchical implementations), M3 uses continuous activation values with
competitive selection, making it particularly suited for theories
emphasizing graded memory strength (for a detailed comparison, see
\citeproc{ref-oberauer2019simple}{Oberauer \& Lewandowsky, 2019}).

\section{Tutorial 1: Simple Span M3}\label{tutorial-1-simple-span-m3}

This tutorial walks through the complete M3 workflow using data from
Oberauer (\citeproc{ref-oberauer2019bindings}{2019}), who measured cued
recall with word pairs. Two experiments (20 participants each) differed
in stimulus set type: Experiment 1 used new word pairs on each trial
(open set), while Experiment 2 reused a small set of words across trials
(closed set). Set size (2, 4, 6, 8 pairs) varied within participants.
The full companion script (\texttt{scripts/tutorial1\_simple\_span.R})
contains all code; we reproduce the key steps here.

\subsection{Data}\label{data}

Each recognition trial presents a probe together with a set of response
options: the correct target, other studied items from the same list, and
not-presented lures. Participants select one item, producing a
categorical response: \emph{correct}, \emph{other} (intrusion from a
different list position), or \emph{NPL} (not-presented lure). The number
of response options per category varies across the 24 recognition
conditions (combinations of set size, list items in the recognition set,
and number of lures).

M3 in \texttt{bmm} expects aggregated data: response frequency counts
per participant and condition, along with the number of response options
per category. Starting from the trial-level data (available at
https://osf.io/vekpd/), we filter to recognition trials, aggregate
response categories, and compute option counts:

\begin{Shaded}
\begin{Highlighting}[]
\CommentTok{\# Read both experiments; create unique participant IDs}
\NormalTok{col\_names }\OtherTok{\textless{}{-}} \FunctionTok{c}\NormalTok{(}\StringTok{"id"}\NormalTok{, }\StringTok{"session"}\NormalTok{, }\StringTok{"block"}\NormalTok{, }\StringTok{"trial"}\NormalTok{, }\StringTok{"setsize"}\NormalTok{,}
               \StringTok{"rsizeList"}\NormalTok{, }\StringTok{"rsizeNPL"}\NormalTok{, }\StringTok{"tested"}\NormalTok{, }\StringTok{"response"}\NormalTok{, }\StringTok{"rcat"}\NormalTok{, }\StringTok{"rt"}\NormalTok{)}

\NormalTok{read\_exp }\OtherTok{\textless{}{-}} \ControlFlowTok{function}\NormalTok{(file, exp\_label) \{}
  \FunctionTok{read.table}\NormalTok{(}\FunctionTok{here}\NormalTok{(}\StringTok{"data"}\NormalTok{, file), }\AttributeTok{header =} \ConstantTok{FALSE}\NormalTok{,}
             \AttributeTok{col.names =}\NormalTok{ col\_names, }\AttributeTok{fill =} \ConstantTok{TRUE}\NormalTok{) }\SpecialCharTok{\%\textgreater{}\%}
    \FunctionTok{drop\_na}\NormalTok{() }\SpecialCharTok{\%\textgreater{}\%}
    \FunctionTok{mutate}\NormalTok{(}\AttributeTok{exp =}\NormalTok{ exp\_label,}
           \AttributeTok{id  =} \FunctionTok{paste}\NormalTok{(id, exp\_label, }\AttributeTok{sep =} \StringTok{"\_"}\NormalTok{))}
\NormalTok{\}}

\NormalTok{data\_raw }\OtherTok{\textless{}{-}} \FunctionTok{bind\_rows}\NormalTok{(}
  \FunctionTok{read\_exp}\NormalTok{(}\StringTok{"Oberauer\_2019\_SimpleSpan\_Exp1.dat"}\NormalTok{, }\StringTok{"openset"}\NormalTok{),}
  \FunctionTok{read\_exp}\NormalTok{(}\StringTok{"Oberauer\_2019\_SimpleSpan\_Exp2.dat"}\NormalTok{, }\StringTok{"closedset"}\NormalTok{)}
\NormalTok{)}
\end{Highlighting}
\end{Shaded}

\begin{Shaded}
\begin{Highlighting}[]
\CommentTok{\# Keep recognition trials only, aggregate to response frequencies}
\NormalTok{data\_agg }\OtherTok{\textless{}{-}}\NormalTok{ data\_raw }\SpecialCharTok{\%\textgreater{}\%}
  \FunctionTok{filter}\NormalTok{(rsizeList }\SpecialCharTok{+}\NormalTok{ rsizeNPL }\SpecialCharTok{\textgreater{}} \DecValTok{0}\NormalTok{) }\SpecialCharTok{\%\textgreater{}\%}
  \FunctionTok{group\_by}\NormalTok{(id, exp, setsize, rsizeList, rsizeNPL) }\SpecialCharTok{\%\textgreater{}\%}
  \FunctionTok{summarise}\NormalTok{(}\AttributeTok{corr  =} \FunctionTok{sum}\NormalTok{(rcat }\SpecialCharTok{==} \DecValTok{1}\NormalTok{),}
            \AttributeTok{other =} \FunctionTok{sum}\NormalTok{(rcat }\SpecialCharTok{==} \DecValTok{2}\NormalTok{),}
            \AttributeTok{npl   =} \FunctionTok{sum}\NormalTok{(rcat }\SpecialCharTok{==} \DecValTok{3}\NormalTok{), }\AttributeTok{.groups =} \StringTok{"drop"}\NormalTok{) }\SpecialCharTok{\%\textgreater{}\%}
  \FunctionTok{mutate}\NormalTok{(}\AttributeTok{n\_corr  =} \DecValTok{1}\NormalTok{,}
         \AttributeTok{n\_other =}\NormalTok{ rsizeList }\SpecialCharTok{{-}} \DecValTok{1}\NormalTok{,  }\CommentTok{\# rsizeList includes the correct item}
         \AttributeTok{n\_npl   =}\NormalTok{ rsizeNPL,}
         \AttributeTok{ss\_lin  =}\NormalTok{ setsize }\SpecialCharTok{{-}} \FunctionTok{mean}\NormalTok{(setsize),}
         \AttributeTok{exp     =} \FunctionTok{factor}\NormalTok{(exp))}
\end{Highlighting}
\end{Shaded}

The resulting dataset has 960 rows (40 participants \(\times\) 24
recognition conditions) with response frequencies (\texttt{corr},
\texttt{other}, \texttt{npl}) and the corresponding option counts
(\texttt{n\_corr}, \texttt{n\_other}, \texttt{n\_npl}). Set size is
mean-centered (\texttt{ss\_lin}) for the linear predictor.

\subsection{Model Specification}\label{model-specification}

Specifying an M3 requires two objects: a model definition
(\texttt{m3()}) and a predictor formula (\texttt{bmf()}). The model
definition declares the response categories, the columns containing
option counts, the model version, and the choice rule. For the simple
span version with the softmax rule:

\begin{Shaded}
\begin{Highlighting}[]
\NormalTok{m3\_model\_ss }\OtherTok{\textless{}{-}} \FunctionTok{m3}\NormalTok{(}
  \AttributeTok{resp\_cats   =} \FunctionTok{c}\NormalTok{(}\StringTok{"corr"}\NormalTok{, }\StringTok{"other"}\NormalTok{, }\StringTok{"npl"}\NormalTok{),}
  \AttributeTok{num\_options =} \FunctionTok{c}\NormalTok{(}\StringTok{"n\_corr"}\NormalTok{, }\StringTok{"n\_other"}\NormalTok{, }\StringTok{"n\_npl"}\NormalTok{),}
  \AttributeTok{version     =} \StringTok{"ss"}\NormalTok{,}
  \AttributeTok{choice\_rule =} \StringTok{"softmax"}
\NormalTok{)}
\end{Highlighting}
\end{Shaded}

The predictor formula specifies how experimental conditions map onto the
activation parameters \(a\) and \(c\). Because the two experiments are
between-subjects, we use cell means coding (\texttt{0\ +\ exp}) so that
each experiment receives its own intercept rather than a difference from
a reference level. Set size enters as a linear predictor
(\texttt{exp:ss\_lin}), and random effects allow both intercepts and
slopes to vary across participants, with separate variance components
per experiment (\texttt{gr(id,\ by\ =\ exp)}):

\begin{Shaded}
\begin{Highlighting}[]
\NormalTok{m3\_formula }\OtherTok{\textless{}{-}} \FunctionTok{bmf}\NormalTok{(}
\NormalTok{  c }\SpecialCharTok{\textasciitilde{}} \DecValTok{0} \SpecialCharTok{+}\NormalTok{ exp }\SpecialCharTok{+}\NormalTok{ exp}\SpecialCharTok{:}\NormalTok{ss\_lin }\SpecialCharTok{+}\NormalTok{ (}\DecValTok{1} \SpecialCharTok{+}\NormalTok{ ss\_lin }\SpecialCharTok{|} \FunctionTok{gr}\NormalTok{(id, }\AttributeTok{by =}\NormalTok{ exp)),}
\NormalTok{  a }\SpecialCharTok{\textasciitilde{}} \DecValTok{0} \SpecialCharTok{+}\NormalTok{ exp }\SpecialCharTok{+}\NormalTok{ exp}\SpecialCharTok{:}\NormalTok{ss\_lin }\SpecialCharTok{+}\NormalTok{ (}\DecValTok{1} \SpecialCharTok{+}\NormalTok{ ss\_lin }\SpecialCharTok{|} \FunctionTok{gr}\NormalTok{(id, }\AttributeTok{by =}\NormalTok{ exp))}
\NormalTok{)}
\end{Highlighting}
\end{Shaded}

Under the softmax rule, \(a\) and \(c\) are estimated on the identity
link (native activation scale). The default priors can be inspected with
\texttt{default\_prior(m3\_formula,\ m3\_model\_ss,\ data\ =\ data\_agg)}.
We use the defaults throughout this tutorial; Tutorial 2 shows how to
set custom priors when needed.

\subsection{Model Fitting}\label{model-fitting}

The \texttt{bmm()} function combines the model definition and formula,
passes them to \texttt{brms} and Stan, and returns a fitted model
object. All models in this tutorial were estimated with 4 chains of
2,000 iterations each (1,000 warmup), using \texttt{cmdstanr} as the
backend. We set \texttt{sample\_prior\ =\ "yes"} so that prior samples
are available for later hypothesis tests. The \texttt{file} argument
caches the fitted model to disk; re-running the script loads the cached
fit rather than re-fitting.

\begin{Shaded}
\begin{Highlighting}[]
\NormalTok{fit\_m3\_ss\_softmax }\OtherTok{\textless{}{-}} \FunctionTok{bmm}\NormalTok{(}
  \AttributeTok{formula      =}\NormalTok{ m3\_formula,}
  \AttributeTok{model        =}\NormalTok{ m3\_model\_ss,}
  \AttributeTok{data         =}\NormalTok{ data\_agg,}
  \AttributeTok{cores        =} \DecValTok{4}\NormalTok{,}
  \AttributeTok{backend      =} \StringTok{"cmdstanr"}\NormalTok{,}
  \AttributeTok{sample\_prior =} \StringTok{"yes"}\NormalTok{,}
  \AttributeTok{save\_pars    =} \FunctionTok{save\_pars}\NormalTok{(}\AttributeTok{all =} \ConstantTok{TRUE}\NormalTok{),}
  \AttributeTok{file         =} \FunctionTok{here}\NormalTok{(}\StringTok{"output"}\NormalTok{, }\StringTok{"fit\_m3\_ss\_softmax"}\NormalTok{)}
\NormalTok{)}
\end{Highlighting}
\end{Shaded}

After fitting, check convergence with
\texttt{summary(fit\_m3\_ss\_softmax)}. All \(\hat{R}\) values
(\citeproc{ref-gelman1992inference}{Gelman \& Rubin, 1992}) should be
below 1.01 and effective sample sizes should be adequate. For all models
in this tutorial, these criteria were met (\(\hat{R} \leq\) 1.013, bulk
ESS \(>\) 460). For guidance on Bayesian convergence diagnostics, see
Gabry et al. (\citeproc{ref-gabry2019visualization}{2019}) and the
\texttt{brms} documentation.

\subsection{Model Evaluation}\label{model-evaluation}

A posterior predictive check compares observed response proportions to
those predicted by the fitted model. We use \texttt{posterior\_epred()}
to obtain expected counts per category for each observation, average
across posterior draws, and aggregate to participant-level proportions
by set size and experiment. \textbf{?@fig-t1-pp-check} shows the result.
Across all conditions and response categories, predicted proportions
fall within one standard error of the observed means.

\begin{Shaded}
\begin{Highlighting}[]
\NormalTok{knitr}\SpecialCharTok{::}\FunctionTok{include\_graphics}\NormalTok{(}\FunctionTok{here}\NormalTok{(}\StringTok{"figures"}\NormalTok{, }\StringTok{"tutorial1\_pp\_check\_softmax.pdf"}\NormalTok{))}
\end{Highlighting}
\end{Shaded}

\subsection{Parameter Estimates}\label{parameter-estimates}

With the softmax choice rule, \(a\) and \(c\) are on the identity link,
so the posterior draws are directly interpretable as activation values.
To compute condition-specific estimates, we reconstruct the linear
predictor from the fixed effects: for each set size \(s\),
\(\hat{a}_{\text{exp}} = \beta_{a,\text{exp}} + \beta_{a,\text{exp:ss\_lin}} \cdot (s - \bar{s})\),
and analogously for \(c\). The companion script extracts draws with
\texttt{as\_draws\_df()} and computes 95\% highest-density intervals
{[}HDIs; McElreath (\citeproc{ref-mcelreath2020statistical}{2020}){]}
via \texttt{tidybayes::mean\_hdci()}.

\textbf{?@fig-t1-param-estimates} shows the resulting estimates. Context
binding (\(c\)) is consistently higher than item memory (\(a\)), which
suggests that the primary source of correct responding is
context-specific retrieval rather than general familiarity. Both
parameters decrease with set size, but the decline is steeper for \(c\),
consistent with binding being more sensitive to memory load. The two
experiments show similar patterns, though the closed set tends to
produce slightly higher item activation, likely because the same words
recur across trials and accumulate familiarity.

\begin{Shaded}
\begin{Highlighting}[]
\NormalTok{knitr}\SpecialCharTok{::}\FunctionTok{include\_graphics}\NormalTok{(}\FunctionTok{here}\NormalTok{(}\StringTok{"figures"}\NormalTok{, }\StringTok{"tutorial1\_param\_estimates\_softmax.pdf"}\NormalTok{))}
\end{Highlighting}
\end{Shaded}

\subsection{Choice Rule Comparison}\label{choice-rule-comparison}

To compare the two choice rules, we fit a second model with
\texttt{choice\_rule\ =\ "simple"} using the same formula and data.
Under the simple rule, \(a\) and \(c\) are estimated on the log link, so
the raw coefficients must be exponentiated to obtain activation values.

The posterior predictive checks for both models are qualitatively
similar, but the softmax rule tends to capture the tails of the response
distribution more accurately, particularly for the NPL category at
larger set sizes. The simple rule slightly overpredicts other-list
intrusions in some conditions. These figures are available in the
companion script.

For a formal comparison, we compute the marginal likelihood of each
model via bridge sampling (\citeproc{ref-gronau2017bridge}{Gronau et
al., 2017}) and derive a Bayes factor. This requires
\texttt{save\_pars(all\ =\ TRUE)} during fitting (set in the
\texttt{bmm()} call above):

\begin{Shaded}
\begin{Highlighting}[]
\NormalTok{bridge\_softmax }\OtherTok{\textless{}{-}} \FunctionTok{bridge\_sampler}\NormalTok{(fit\_m3\_ss\_softmax, }\AttributeTok{repetitions =} \DecValTok{10}\NormalTok{)}
\NormalTok{bridge\_simple  }\OtherTok{\textless{}{-}} \FunctionTok{bridge\_sampler}\NormalTok{(fit\_m3\_ss\_simple,  }\AttributeTok{repetitions =} \DecValTok{10}\NormalTok{)}
\FunctionTok{bf}\NormalTok{(bridge\_softmax, bridge\_simple)}
\end{Highlighting}
\end{Shaded}

The Bayes factor strongly favors the softmax rule
(\emph{BF}\textsubscript{10} = \(2.7 \times 10^{7}\),
\(3.2 \times 10^{7}\), \(3.1 \times 10^{7}\), \(2.0 \times 10^{7}\),
\(2.6 \times 10^{7}\), \(2.9 \times 10^{7}\), \(2.8 \times 10^{7}\),
\(3.9 \times 10^{7}\), \(1.2 \times 10^{7}\), \(2.8 \times 10^{7}\)).
This is consistent with recent comparisons by Oberauer
(\citeproc{ref-oberauer2023measurement}{2023}) and Li et al.
(\citeproc{ref-li2026filtering}{in press}), who found softmax to
outperform the simple rule across different working memory tasks. We use
the softmax rule as the default in the remaining tutorials.

\subsection{Hypothesis Testing}\label{hypothesis-testing}

Because we fitted the model with \texttt{sample\_prior\ =\ "yes"}, we
can use \texttt{brms::hypothesis()} to test directed hypotheses and
obtain evidence ratios (Savage-Dickey density ratios comparing posterior
to prior density at the hypothesis boundary). For example, to test
whether context binding exceeds item memory within the open-set
experiment:

\begin{Shaded}
\begin{Highlighting}[]
\FunctionTok{hypothesis}\NormalTok{(fit\_m3\_ss\_softmax, }\StringTok{"c\_expopenset {-} a\_expopenset \textgreater{} 0"}\NormalTok{)}
\end{Highlighting}
\end{Shaded}

This returns an evidence ratio quantifying how much the data shift
belief toward \(c > a\) relative to the prior (ER = 58.7 for the
open-set experiment). We can similarly test whether set size reduces
\(c\) (averaged across experiments) and whether set size effects differ
between experiments:

\begin{Shaded}
\begin{Highlighting}[]
\CommentTok{\# Set size effect on context binding (averaged across experiments)}
\FunctionTok{hypothesis}\NormalTok{(fit\_m3\_ss\_softmax,}
  \StringTok{"(c\_expopenset:ss\_lin + c\_expclosedset:ss\_lin) / 2 \textless{} 0"}\NormalTok{)}

\CommentTok{\# Experiment difference in set size effect on c}
\FunctionTok{hypothesis}\NormalTok{(fit\_m3\_ss\_softmax,}
  \StringTok{"c\_expopenset:ss\_lin {-} c\_expclosedset:ss\_lin = 0"}\NormalTok{)}
\end{Highlighting}
\end{Shaded}

The results confirm that \(c > a\) in both experiments, that set size
reduces both \(a\) and \(c\), and that the set size effect on \(c\) is
steeper for the closed-set experiment. The companion script contains the
full set of six hypothesis tests. For other examples of M3-based
hypothesis testing with simple span data, see Loaiza and Souza
(\citeproc{ref-loaiza2024active}{2024}); for an application of the cs
version to WM load and intelligence, see Schubert et al.
(\citeproc{ref-schubert2023working}{2023}).

\section{Tutorial 2: Complex Span M3}\label{tutorial-2-complex-span-m3}

This tutorial extends the workflow to tasks where distractors appear
alongside memoranda. We use data from Li et al.
(\citeproc{ref-li2026filtering}{in press}), who developed a complex-span
filtering task to test whether processed information can be kept out of
working memory. The companion script is
\texttt{scripts/tutorial2\_complex\_span.R}; data preparation is in
\texttt{scripts/prepare\_Li2026\_data.R}.

\subsection{What makes complex span
different}\label{what-makes-complex-span-different}

The simple span model has three response categories. When distractors
are present, two more categories are needed: the distractor paired with
the tested target (\emph{paired distractor}) and distractors from other
positions (\emph{other distractor}). The complex span (cs) model extends
the activation formulas with a filtering parameter \(f\) that scales the
activation of distractor categories:

\[
\begin{aligned}
\text{act}_{\text{correct}} &= b + a + c \\
\text{act}_{\text{paired dist.}} &= b + f \cdot a + f \cdot c \\
\text{act}_{\text{other}} &= b + a \\
\text{act}_{\text{other dist.}} &= b + f \cdot a \\
\text{act}_{\text{NPL}} &= b
\end{aligned}
\]

The parameter \(f\) ranges from 0 to 1 (estimated on the logit scale).
When \(f = 0\), distractors receive only baseline activation (perfect
filtering). When \(f = 1\), distractors are activated as strongly as
memoranda (no filtering).

A critical implementation detail: the \texttt{cs} version assigns
activation formulas by position, so \texttt{resp\_cats} must follow the
order \texttt{corr,\ distc,\ other,\ disto,\ npl}. Misordering silently
assigns the wrong formula to each category.

\subsection{Data}\label{data-1}

Li et al. (\citeproc{ref-li2026filtering}{in press}) presented
participants with image pairs and manipulated whether targets and
distractors were identified before processing (pre-cue), after
processing (retro-cue), or whether no distractors were present
(control). We use Experiment 1 (45 participants after exclusions,
within-subjects). The data are pre-aggregated into response frequencies
across the five categories; see the companion preparation script for
classification details.

\begin{Shaded}
\begin{Highlighting}[]
\NormalTok{data\_agg }\OtherTok{\textless{}{-}} \FunctionTok{read\_csv}\NormalTok{(}\FunctionTok{here}\NormalTok{(}\StringTok{"data"}\NormalTok{, }\StringTok{"Li\_2026\_ComplexSpan\_Exp1\_agg.csv"}\NormalTok{),}
                     \AttributeTok{show\_col\_types =} \ConstantTok{FALSE}\NormalTok{) }\SpecialCharTok{\%\textgreater{}\%}
  \FunctionTok{mutate}\NormalTok{(}\AttributeTok{condition =} \FunctionTok{factor}\NormalTok{(condition, }\AttributeTok{levels =} \FunctionTok{c}\NormalTok{(}\StringTok{"control"}\NormalTok{, }\StringTok{"pre"}\NormalTok{, }\StringTok{"retro"}\NormalTok{)))}
\end{Highlighting}
\end{Shaded}

\subsection{Model specification}\label{model-specification-1}

The model definition follows the same pattern as Tutorial 1:

\begin{Shaded}
\begin{Highlighting}[]
\NormalTok{m3\_model\_cs }\OtherTok{\textless{}{-}} \FunctionTok{m3}\NormalTok{(}
  \AttributeTok{resp\_cats   =} \FunctionTok{c}\NormalTok{(}\StringTok{"corr"}\NormalTok{, }\StringTok{"distc"}\NormalTok{, }\StringTok{"other"}\NormalTok{, }\StringTok{"disto"}\NormalTok{, }\StringTok{"npl"}\NormalTok{),}
  \AttributeTok{num\_options =} \FunctionTok{c}\NormalTok{(}\StringTok{"n\_corr"}\NormalTok{, }\StringTok{"n\_distc"}\NormalTok{, }\StringTok{"n\_other"}\NormalTok{, }\StringTok{"n\_disto"}\NormalTok{, }\StringTok{"n\_npl"}\NormalTok{),}
  \AttributeTok{version     =} \StringTok{"cs"}\NormalTok{,}
  \AttributeTok{choice\_rule =} \StringTok{"softmax"}
\NormalTok{)}
\end{Highlighting}
\end{Shaded}

The predictor formula specifies condition effects on all three
parameters with cell means coding and uncorrelated random effects per
condition:

\begin{Shaded}
\begin{Highlighting}[]
\NormalTok{m3\_formula\_cs }\OtherTok{\textless{}{-}} \FunctionTok{bmf}\NormalTok{(}
\NormalTok{  c }\SpecialCharTok{\textasciitilde{}} \DecValTok{0} \SpecialCharTok{+}\NormalTok{ condition }\SpecialCharTok{+}\NormalTok{ (}\DecValTok{0} \SpecialCharTok{+}\NormalTok{ condition }\SpecialCharTok{||}\NormalTok{ participant),}
\NormalTok{  a }\SpecialCharTok{\textasciitilde{}} \DecValTok{0} \SpecialCharTok{+}\NormalTok{ condition }\SpecialCharTok{+}\NormalTok{ (}\DecValTok{0} \SpecialCharTok{+}\NormalTok{ condition }\SpecialCharTok{||}\NormalTok{ participant),}
\NormalTok{  f }\SpecialCharTok{\textasciitilde{}} \DecValTok{0} \SpecialCharTok{+}\NormalTok{ condition }\SpecialCharTok{+}\NormalTok{ (}\DecValTok{0} \SpecialCharTok{+}\NormalTok{ condition }\SpecialCharTok{||}\NormalTok{ participant)}
\NormalTok{)}
\end{Highlighting}
\end{Shaded}

There is a complication: in the control condition, no distractors are
present (\texttt{n\_distc\ =\ 0}, \texttt{n\_disto\ =\ 0}), so \(f\)
does not influence the likelihood and is not identified. If left free,
the sampler will wander through the prior, and the resulting large
variance can affect other parameters. The solution is to fix \(f\) in
the control condition using a \texttt{constant()} prior on both the
fixed effect and the random effect SD:

\begin{Shaded}
\begin{Highlighting}[]
\NormalTok{priors\_cs }\OtherTok{\textless{}{-}} \FunctionTok{c}\NormalTok{(}
  \FunctionTok{prior}\NormalTok{(}\FunctionTok{constant}\NormalTok{(}\DecValTok{10}\NormalTok{), }\AttributeTok{nlpar =} \StringTok{"f"}\NormalTok{, }\AttributeTok{coef =} \StringTok{"conditioncontrol"}\NormalTok{),}
  \FunctionTok{prior}\NormalTok{(}\FunctionTok{constant}\NormalTok{(}\DecValTok{0}\NormalTok{),  }\AttributeTok{class =} \StringTok{"sd"}\NormalTok{, }\AttributeTok{nlpar =} \StringTok{"f"}\NormalTok{,}
        \AttributeTok{coef =} \StringTok{"conditioncontrol"}\NormalTok{, }\AttributeTok{group =} \StringTok{"participant"}\NormalTok{)}
\NormalTok{)}
\end{Highlighting}
\end{Shaded}

Setting the fixed effect to 10 on the logit scale (effectively
\(f \approx 1\)) and the SD to 0 removes \(f\) in the control condition
from the model entirely.

\subsection{Model fitting and
evaluation}\label{model-fitting-and-evaluation}

Fitting follows the same \texttt{bmm()} call as in Tutorial 1, passing
the custom priors. The model converged without issues (\(\hat{R} \leq\)
1.004, bulk ESS \(>\) 750, no divergent transitions).
\textbf{?@fig-t2-pp-check} shows the posterior predictive check. The
model captures the observed response proportions well across all three
conditions and five categories.

\begin{Shaded}
\begin{Highlighting}[]
\NormalTok{knitr}\SpecialCharTok{::}\FunctionTok{include\_graphics}\NormalTok{(}\FunctionTok{here}\NormalTok{(}\StringTok{"figures"}\NormalTok{, }\StringTok{"tutorial2\_pp\_check.pdf"}\NormalTok{))}
\end{Highlighting}
\end{Shaded}

\subsection{Parameter estimates and interpreting
f}\label{parameter-estimates-and-interpreting-f}

\textbf{?@fig-t2-param-estimates} shows the posterior estimates. Item
memory (\(a\)) and context binding (\(c\)) are on the identity link and
can be read directly. The filtering parameter \(f\) is estimated on the
logit scale; to interpret it as a probability, we apply the inverse
logit transformation (\texttt{plogis()} in R):

\begin{Shaded}
\begin{Highlighting}[]
\NormalTok{draws }\OtherTok{\textless{}{-}} \FunctionTok{as\_draws\_df}\NormalTok{(fit\_m3\_cs) }\SpecialCharTok{\%\textgreater{}\%} \FunctionTok{as\_tibble}\NormalTok{()}
\NormalTok{f\_pre   }\OtherTok{\textless{}{-}} \FunctionTok{plogis}\NormalTok{(draws}\SpecialCharTok{$}\NormalTok{b\_f\_conditionpre)}
\NormalTok{f\_retro }\OtherTok{\textless{}{-}} \FunctionTok{plogis}\NormalTok{(draws}\SpecialCharTok{$}\NormalTok{b\_f\_conditionretro)}
\end{Highlighting}
\end{Shaded}

\begin{Shaded}
\begin{Highlighting}[]
\NormalTok{knitr}\SpecialCharTok{::}\FunctionTok{include\_graphics}\NormalTok{(}\FunctionTok{here}\NormalTok{(}\StringTok{"figures"}\NormalTok{, }\StringTok{"tutorial2\_param\_estimates.pdf"}\NormalTok{))}
\end{Highlighting}
\end{Shaded}

On the probability scale, \(f\) is close to 1 in both conditions
(pre-cue: \(M\) = 0.66, 95\% HDI {[}0.61, 0.70{]}; retro-cue: \(M\) =
0.63, 95\% HDI {[}0.59, 0.66{]}), indicating little to no filtering.
This matches the conclusion of Li et al.
(\citeproc{ref-li2026filtering}{in press}): once distractors are
processed, they are encoded into working memory regardless of whether
participants knew in advance which items were distractors.

\subsection{Hypothesis tests}\label{hypothesis-tests}

We can test condition differences formally using
\texttt{brms::hypothesis()}:

\begin{Shaded}
\begin{Highlighting}[]
\CommentTok{\# Does f differ between pre{-}cue and retro{-}cue?}
\FunctionTok{hypothesis}\NormalTok{(fit\_m3\_cs, }\StringTok{"f\_conditionpre {-} f\_conditionretro = 0"}\NormalTok{)}

\CommentTok{\# Do a and c differ between control and distractor conditions?}
\FunctionTok{hypothesis}\NormalTok{(fit\_m3\_cs,}
  \StringTok{"a\_conditioncontrol {-} (a\_conditionpre + a\_conditionretro) / 2 = 0"}\NormalTok{)}
\end{Highlighting}
\end{Shaded}

The evidence ratio for the \(f\) comparison is close to 1 (ER = 10.3),
providing no evidence for a difference in filtering between cue
conditions. Item memory and context binding are slightly lower in the
distractor conditions than in the control, in line with a general cost
of distractor processing on memory strength.

\subsection{Comparison with the published
model}\label{comparison-with-the-published-model}

The standard \texttt{cs} model constrains distractor attenuation to be
equal for item memory and context binding: both are scaled by the single
parameter \(f\). The original analysis by Li et al.
(\citeproc{ref-li2026filtering}{in press}) used a custom model with two
separate ratio parameters, \(r_a\) and \(r_c\), allowing item and
binding activation to transfer differently to distractors. The
\texttt{cs} model implicitly assumes \(r_a = r_c = f\).

When this equal-attenuation constraint is too restrictive for the
research question, a custom model is needed. Tutorial 3 demonstrates how
to define such models from scratch.

\section{Tutorial 3: Custom M3 --- Parameter
Recovery}\label{tutorial-3-custom-m3-parameter-recovery}

The previous tutorials used pre-built model versions (\texttt{ss} and
\texttt{cs}). When a task design does not fit one of these templates,
\texttt{bmm} allows defining a fully custom M3 with user-specified
response categories, activation formulas, and link functions. This
tutorial demonstrates the custom version through a parameter recovery
simulation: we define a custom model, simulate data with known
parameters, fit the model, and check whether the generating values are
recovered. The walkthrough follows
\texttt{scripts/tutorial3\_parameter\_recovery\_simple.R}; design-grid
results come from \texttt{scripts/tutorial3\_parameter\_recovery.R}.

\subsection{Why parameter recovery}\label{why-parameter-recovery}

Before applying a new or modified M3 to empirical data, the model should
be able to recover the parameters that generated the data. If recovery
fails under ideal conditions (known generating values, no model
misspecification), the problem lies in the model or the design, not in
the data. Recovery simulations are especially informative for custom
models where the number of parameters and the functional form of the
activation equations are chosen by the researcher Wilson \& Collins
(\citeproc{ref-wilson2019ten}{2019}).

\subsection{Defining a custom M3}\label{defining-a-custom-m3}

As a running example, we use a memory updating task from Oberauer and
Lewandowsky (\citeproc{ref-oberauer2019simple}{2019}; applied with a
custom M3 by \citeproc{ref-li2025updating}{Li et al., 2025}).
Participants memorize items at five positions, then each position is
updated once (the old item is replaced). At test, one position is probed
and the participant selects from a recognition set containing: the
current target (1 option), other current items (4), the old item from
the tested position (1), old items from other positions (4), and
not-presented lures (5). This produces five response categories with the
following activation formulas:

\[
\begin{aligned}
\text{act}_{\text{correct}}  &= b + (1 + e \cdot t_e) \cdot (a + c) \\
\text{act}_{\text{other}}    &= b + (1 + e \cdot t_e) \cdot a \\
\text{act}_{\text{oldinpos}} &= b + \exp(-r \cdot t_r) \cdot d \cdot (1 + e \cdot t_e) \cdot (a + c) \\
\text{act}_{\text{otherold}} &= b + \exp(-r \cdot t_r) \cdot d \cdot (1 + e \cdot t_e) \cdot a \\
\text{act}_{\text{NPL}}      &= b
\end{aligned}
\]

Beyond the familiar \(a\) (item memory) and \(c\) (context binding),
three new parameters appear: \(d\) captures the residual activation of
replaced (outdated) items, \(e\) scales the benefit of extended encoding
time \(t_e\), and \(r\) governs how quickly old-item activation decays
with removal time \(t_r\). The design variables \(t_e\) and \(t_r\) are
experimentally manipulated (each at 0.25, 0.5, and 1.0 s), producing a
\(3 \times 3\) within-subjects design with 9 conditions.

The model definition in \texttt{bmm} specifies the response categories,
option counts, and link functions. Link functions are chosen to match
each parameter's natural scale: identity for \(a\) and \(c\) (unbounded
activation), logit for \(d\) (bounded between 0 and 1), and log for
\(e\) and \(r\) (positive values):

\begin{Shaded}
\begin{Highlighting}[]
\NormalTok{m3\_model\_custom }\OtherTok{\textless{}{-}} \FunctionTok{m3}\NormalTok{(}
  \AttributeTok{resp\_cats =} \FunctionTok{c}\NormalTok{(}\StringTok{"correct"}\NormalTok{, }\StringTok{"other"}\NormalTok{, }\StringTok{"oldinpos"}\NormalTok{, }\StringTok{"otherold"}\NormalTok{, }\StringTok{"npl"}\NormalTok{),}
  \AttributeTok{num\_options =} \FunctionTok{c}\NormalTok{(}
    \StringTok{"n\_correct"} \OtherTok{=} \DecValTok{1}\NormalTok{, }\StringTok{"n\_other"} \OtherTok{=} \DecValTok{4}\NormalTok{, }\StringTok{"n\_oldinpos"} \OtherTok{=} \DecValTok{1}\NormalTok{,}
    \StringTok{"n\_otherold"} \OtherTok{=} \DecValTok{4}\NormalTok{, }\StringTok{"n\_npl"} \OtherTok{=} \DecValTok{5}
\NormalTok{  ),}
  \AttributeTok{version     =} \StringTok{"custom"}\NormalTok{,}
  \AttributeTok{choice\_rule =} \StringTok{"softmax"}\NormalTok{,}
  \AttributeTok{links =} \FunctionTok{list}\NormalTok{(}
    \AttributeTok{a =} \StringTok{"identity"}\NormalTok{, }\AttributeTok{c =} \StringTok{"identity"}\NormalTok{,}
    \AttributeTok{d =} \StringTok{"logit"}\NormalTok{, }\AttributeTok{e =} \StringTok{"log"}\NormalTok{, }\AttributeTok{r =} \StringTok{"log"}
\NormalTok{  )}
\NormalTok{)}
\end{Highlighting}
\end{Shaded}

The activation formulas are specified separately using \texttt{bmf()}.
Note that \texttt{te} and \texttt{tr} are data variables (not estimated
parameters); \texttt{bmm} parses these expressions and passes them to
\texttt{brms} as non-linear formulas:

\begin{Shaded}
\begin{Highlighting}[]
\NormalTok{m3\_act\_funs }\OtherTok{\textless{}{-}} \FunctionTok{bmf}\NormalTok{(}
\NormalTok{  correct  }\SpecialCharTok{\textasciitilde{}}\NormalTok{ (}\DecValTok{1} \SpecialCharTok{+}\NormalTok{ e }\SpecialCharTok{*}\NormalTok{ te) }\SpecialCharTok{*}\NormalTok{ (a }\SpecialCharTok{+}\NormalTok{ c) }\SpecialCharTok{+}\NormalTok{ b,}
\NormalTok{  other    }\SpecialCharTok{\textasciitilde{}}\NormalTok{ (}\DecValTok{1} \SpecialCharTok{+}\NormalTok{ e }\SpecialCharTok{*}\NormalTok{ te) }\SpecialCharTok{*}\NormalTok{ a }\SpecialCharTok{+}\NormalTok{ b,}
\NormalTok{  oldinpos }\SpecialCharTok{\textasciitilde{}} \FunctionTok{exp}\NormalTok{(}\SpecialCharTok{{-}}\NormalTok{r }\SpecialCharTok{*}\NormalTok{ tr) }\SpecialCharTok{*}\NormalTok{ d }\SpecialCharTok{*}\NormalTok{ (}\DecValTok{1} \SpecialCharTok{+}\NormalTok{ e }\SpecialCharTok{*}\NormalTok{ te) }\SpecialCharTok{*}\NormalTok{ (a }\SpecialCharTok{+}\NormalTok{ c) }\SpecialCharTok{+}\NormalTok{ b,}
\NormalTok{  otherold }\SpecialCharTok{\textasciitilde{}} \FunctionTok{exp}\NormalTok{(}\SpecialCharTok{{-}}\NormalTok{r }\SpecialCharTok{*}\NormalTok{ tr) }\SpecialCharTok{*}\NormalTok{ d }\SpecialCharTok{*}\NormalTok{ (}\DecValTok{1} \SpecialCharTok{+}\NormalTok{ e }\SpecialCharTok{*}\NormalTok{ te) }\SpecialCharTok{*}\NormalTok{ a }\SpecialCharTok{+}\NormalTok{ b,}
\NormalTok{  npl      }\SpecialCharTok{\textasciitilde{}}\NormalTok{ b}
\NormalTok{)}
\end{Highlighting}
\end{Shaded}

\subsection{Choosing generating parameter
values}\label{choosing-generating-parameter-values}

For the recovery simulation, we need group-level means and
between-person standard deviations for each parameter. These are
specified on the link scale (the scale on which the sampler operates),
not the native scale. Table~\ref{tbl-gen-pars} lists the generating
values. The native range column shows the approximate span of individual
differences (\(\pm 2\) SD around the mean, transformed to native scale
via the inverse link functions).

\begin{table}

{\caption{{Generating parameter values for the recovery simulation. All
values are on the link scale.}{\label{tbl-gen-pars}}}
\vspace{-20pt}}

\begin{longtable}[]{@{}
  >{\raggedright\arraybackslash}p{(\linewidth - 8\tabcolsep) * \real{0.1618}}
  >{\raggedright\arraybackslash}p{(\linewidth - 8\tabcolsep) * \real{0.0882}}
  >{\raggedright\arraybackslash}p{(\linewidth - 8\tabcolsep) * \real{0.1912}}
  >{\raggedright\arraybackslash}p{(\linewidth - 8\tabcolsep) * \real{0.1618}}
  >{\raggedright\arraybackslash}p{(\linewidth - 8\tabcolsep) * \real{0.3971}}@{}}
\toprule\noalign{}
\begin{minipage}[b]{\linewidth}\raggedright
Parameter
\end{minipage} & \begin{minipage}[b]{\linewidth}\raggedright
Link
\end{minipage} & \begin{minipage}[b]{\linewidth}\raggedright
Mean (link)
\end{minipage} & \begin{minipage}[b]{\linewidth}\raggedright
SD (link)
\end{minipage} & \begin{minipage}[b]{\linewidth}\raggedright
Native range (\(\pm 2\) SD)
\end{minipage} \\
\midrule\noalign{}
\endhead
\bottomrule\noalign{}
\endlastfoot
\(c\) & identity & 2.0 & 0.50 & 1.0 -- 3.0 \\
\(a\) & identity & 1.0 & 0.50 & 0.0 -- 2.0 \\
\(d\) & logit & 0.85 & 0.65 & 0.38 -- 0.89 \\
\(e\) & log & 0.00 & 0.35 & 0.50 -- 2.01 \\
\(r\) & log & \$-\$0.69 & 0.38 & 0.24 -- 1.04 \\
\end{longtable}

\end{table}

Individual parameters are drawn from
\(\mathcal{N}(\mu_\text{link}, \sigma_\text{link})\) independently for
each person and parameter. We simulate \(N = 50\) participants, each
completing 25 trials per condition (225 trials total). With set size 5,
this corresponds to roughly one hour of testing.

\subsection{Simulating data with rm3()}\label{simulating-data-with-rm3}

The \texttt{rm3()} function generates multinomial response counts from a
custom M3 specification. It takes a vector of parameter values (on the
native scale), the model definition, the activation formulas, and the
design variables as inputs. Here is a call for a single participant in
one condition:

\begin{Shaded}
\begin{Highlighting}[]
\NormalTok{demo\_data }\OtherTok{\textless{}{-}} \FunctionTok{rm3}\NormalTok{(}
  \AttributeTok{n =} \DecValTok{1}\NormalTok{, }\AttributeTok{size =} \DecValTok{25}\NormalTok{,}
  \AttributeTok{pars =} \FunctionTok{c}\NormalTok{(}\AttributeTok{a =} \FloatTok{1.2}\NormalTok{, }\AttributeTok{c =} \FloatTok{2.3}\NormalTok{, }\AttributeTok{d =} \FloatTok{0.7}\NormalTok{, }\AttributeTok{e =} \FloatTok{1.1}\NormalTok{, }\AttributeTok{r =} \FloatTok{0.5}\NormalTok{),}
  \AttributeTok{m3\_model =}\NormalTok{ m3\_model\_custom,}
  \AttributeTok{act\_funs =}\NormalTok{ m3\_act\_funs,}
  \AttributeTok{te =} \FloatTok{0.25}\NormalTok{, }\AttributeTok{tr =} \FloatTok{0.25}
\NormalTok{)}
\end{Highlighting}
\end{Shaded}

This returns a one-row data frame with response counts across the five
categories. The full simulation loops over all 50 participants and 9
conditions, generating 450 rows of simulated data. The companion script
contains the complete loop and data inspection code.

\subsection{Model fitting}\label{model-fitting-1}

We fit the simulated data with the same \texttt{bmm()} workflow used in
the previous tutorials, now passing the custom model object:

\begin{Shaded}
\begin{Highlighting}[]
\CommentTok{\# Prediction formula: intercept + random intercept per participant}
\NormalTok{m3\_pred\_funs }\OtherTok{\textless{}{-}} \FunctionTok{bmf}\NormalTok{(}
\NormalTok{  correct  }\SpecialCharTok{\textasciitilde{}}\NormalTok{ (}\DecValTok{1} \SpecialCharTok{+}\NormalTok{ e }\SpecialCharTok{*}\NormalTok{ te) }\SpecialCharTok{*}\NormalTok{ (a }\SpecialCharTok{+}\NormalTok{ c) }\SpecialCharTok{+}\NormalTok{ b,}
\NormalTok{  other    }\SpecialCharTok{\textasciitilde{}}\NormalTok{ (}\DecValTok{1} \SpecialCharTok{+}\NormalTok{ e }\SpecialCharTok{*}\NormalTok{ te) }\SpecialCharTok{*}\NormalTok{ a }\SpecialCharTok{+}\NormalTok{ b,}
\NormalTok{  oldinpos }\SpecialCharTok{\textasciitilde{}} \FunctionTok{exp}\NormalTok{(}\SpecialCharTok{{-}}\NormalTok{r }\SpecialCharTok{*}\NormalTok{ tr) }\SpecialCharTok{*}\NormalTok{ d }\SpecialCharTok{*}\NormalTok{ (}\DecValTok{1} \SpecialCharTok{+}\NormalTok{ e }\SpecialCharTok{*}\NormalTok{ te) }\SpecialCharTok{*}\NormalTok{ (a }\SpecialCharTok{+}\NormalTok{ c) }\SpecialCharTok{+}\NormalTok{ b,}
\NormalTok{  otherold }\SpecialCharTok{\textasciitilde{}} \FunctionTok{exp}\NormalTok{(}\SpecialCharTok{{-}}\NormalTok{r }\SpecialCharTok{*}\NormalTok{ tr) }\SpecialCharTok{*}\NormalTok{ d }\SpecialCharTok{*}\NormalTok{ (}\DecValTok{1} \SpecialCharTok{+}\NormalTok{ e }\SpecialCharTok{*}\NormalTok{ te) }\SpecialCharTok{*}\NormalTok{ a }\SpecialCharTok{+}\NormalTok{ b,}
\NormalTok{  npl      }\SpecialCharTok{\textasciitilde{}}\NormalTok{ b,}
\NormalTok{  a }\SpecialCharTok{\textasciitilde{}} \DecValTok{1} \SpecialCharTok{+}\NormalTok{ (}\DecValTok{1} \SpecialCharTok{||}\NormalTok{ id), c }\SpecialCharTok{\textasciitilde{}} \DecValTok{1} \SpecialCharTok{+}\NormalTok{ (}\DecValTok{1} \SpecialCharTok{||}\NormalTok{ id),}
\NormalTok{  d }\SpecialCharTok{\textasciitilde{}} \DecValTok{1} \SpecialCharTok{+}\NormalTok{ (}\DecValTok{1} \SpecialCharTok{||}\NormalTok{ id), e }\SpecialCharTok{\textasciitilde{}} \DecValTok{1} \SpecialCharTok{+}\NormalTok{ (}\DecValTok{1} \SpecialCharTok{||}\NormalTok{ id), r }\SpecialCharTok{\textasciitilde{}} \DecValTok{1} \SpecialCharTok{+}\NormalTok{ (}\DecValTok{1} \SpecialCharTok{||}\NormalTok{ id)}
\NormalTok{)}

\NormalTok{fit }\OtherTok{\textless{}{-}} \FunctionTok{bmm}\NormalTok{(}
  \AttributeTok{formula =}\NormalTok{ m3\_pred\_funs,}
  \AttributeTok{model   =}\NormalTok{ m3\_model\_custom,}
  \AttributeTok{data    =}\NormalTok{ sim\_data,}
  \AttributeTok{cores   =} \DecValTok{4}\NormalTok{,}
  \AttributeTok{backend =} \StringTok{"cmdstanr"}\NormalTok{,}
  \AttributeTok{file    =} \FunctionTok{here}\NormalTok{(}\StringTok{"output"}\NormalTok{, }\StringTok{"fit\_m3\_custom\_simple\_N50\_T25"}\NormalTok{)}
\NormalTok{)}
\end{Highlighting}
\end{Shaded}

Note that the prediction formula repeats the activation equations and
adds regression formulas for each parameter (here, an intercept plus a
random intercept per participant). Fitting a custom model with 5
parameters and 50 participants takes roughly 10--30 minutes depending on
hardware.

\subsection{Recovery evaluation}\label{recovery-evaluation}

\subsubsection{Group-level recovery}\label{group-level-recovery}

We first compare the recovered fixed effects to the true generating
means. For each parameter, we extract the posterior median and 95\% HDI
from \texttt{fixef()} and check two things: (1) bias (is the posterior
median close to the true value?) and (2) coverage (does the 95\% HDI
contain the true value?). We also recover the between-person SDs from
\texttt{VarCorr()} and compare them to the generating SDs.

\subsubsection{Individual-level
recovery}\label{individual-level-recovery}

Individual-level recovery asks a finer question: for each participant,
does the recovered estimate (fixed + random effect, extracted via
\texttt{coef()}) match their true generating value?
\textbf{?@fig-t3-indiv-scatter} shows scatter plots of true versus
recovered parameters for the single simulation cell (\(N = 50\), 25
trials per condition). Points close to the diagonal indicate good
recovery; the Pearson correlation (\(r\)) is annotated in each panel.

\begin{Shaded}
\begin{Highlighting}[]
\NormalTok{knitr}\SpecialCharTok{::}\FunctionTok{include\_graphics}\NormalTok{(}\FunctionTok{here}\NormalTok{(}\StringTok{"figures"}\NormalTok{, }\StringTok{"tutorial3\_simple\_indiv\_scatter.pdf"}\NormalTok{))}
\end{Highlighting}
\end{Shaded}

Parameters \(a\), \(c\), and \(d\) show strong individual-level recovery
(\(r\) = 0.84, 0.91, and 0.74, respectively). The parameters \(e\) and
\(r\) are harder to recover (\(r\) = 0.70 and 0.43) because they
interact multiplicatively with the design variables \(t_e\) and \(t_r\),
which partially confounds them.

\subsection{Practical guidance}\label{practical-guidance}

The simplified script demonstrates the workflow for a single design
cell. To understand how recovery depends on sample size and the number
of trials per condition, the full simulation script
(\texttt{tutorial3\_parameter\_recovery.R}) varies both factors across a
\(3 \times 3\) grid: \(N \in \{25, 50, 100\}\) and trials per condition
\(\in \{5, 10, 25\}\).

\textbf{?@fig-t3-group-recovery} shows group-level recovery across the
design grid. Recovered estimates are close to the true values in most
cells, with HDIs narrowing as \(N\) increases. With only 5 trials per
condition, some parameters (particularly \(e\) and \(r\)) show
noticeable bias, but this improves substantially at 10 and 25 trials.

\begin{Shaded}
\begin{Highlighting}[]
\NormalTok{knitr}\SpecialCharTok{::}\FunctionTok{include\_graphics}\NormalTok{(}\FunctionTok{here}\NormalTok{(}\StringTok{"figures"}\NormalTok{, }\StringTok{"tutorial3\_group\_recovery.pdf"}\NormalTok{))}
\end{Highlighting}
\end{Shaded}

\textbf{?@fig-t3-heatmap} shows individual-level recovery (Pearson
correlation between true and recovered values) for each parameter across
the same design grid. Individual-level recovery is more demanding than
group-level recovery: \(e\) and \(r\) require both a reasonable sample
size and a sufficient number of trials per condition to reach
correlations above .7.

\begin{Shaded}
\begin{Highlighting}[]
\NormalTok{knitr}\SpecialCharTok{::}\FunctionTok{include\_graphics}\NormalTok{(}\FunctionTok{here}\NormalTok{(}\StringTok{"figures"}\NormalTok{, }\StringTok{"tutorial3\_recovery\_heatmap.pdf"}\NormalTok{))}
\end{Highlighting}
\end{Shaded}

These results are based on a single replication per design cell. A
proper simulation study should replicate each cell many times to account
for Monte Carlo variability Talts et al.
(\citeproc{ref-talts2018validating}{2018}). Recovery quality also
depends on the model structure itself. Parameters that appear in
interaction terms or that are only weakly identified by the design may
require more data. Running a recovery check before collecting empirical
data helps determine whether the planned sample size and number of
trials are sufficient.

\section{Summary and Additional
Resources}\label{summary-and-additional-resources}

This tutorial covered three progressively complex applications of the M3
framework using the \texttt{bmm} R package. The simple span example
(Tutorial 1) introduced the full workflow: specification, fitting,
evaluation, and hypothesis testing. It also compared the two choice
rules. The complex span example (Tutorial 2) added distractor
categories, the filtering parameter \(f\), and the handling of
non-identified parameters. The custom model example (Tutorial 3) showed
how to define an M3 from scratch for a memory updating task and how to
validate it through parameter recovery.

A few practical recommendations. If the task produces three response
categories (correct, other-list intrusion, NPL), the simple span version
(\texttt{ss}) is the natural starting point. Tasks with distractors that
are the same material type as memoranda fit the complex span version
(\texttt{cs}). For other designs, the custom version gives full control
over activation formulas, parameters, and link functions. In all cases,
the softmax choice rule is the recommended default based on current
evidence. When defining a custom model, running a parameter recovery
simulation before data collection reveals whether the planned design can
distinguish the parameters of interest and whether the model is
identified.

\subsection{Additional resources}\label{additional-resources}

\begin{itemize}
\tightlist
\item
  \texttt{bmm} package documentation and M3 vignette:
  \url{https://venpopov.github.io/bmm/articles/bmm_m3.html}
\item
  Full companion scripts for all three tutorials: {[}repository link{]}
\item
  M3 framework paper: Oberauer and Lewandowsky
  (\citeproc{ref-oberauer2019simple}{2019})
\item
  Systematic parameter recovery for M3: Göttmann et al.
  (\citeproc{ref-goettmann2025recovery}{2025})
\item
  Companion \texttt{bmm} tutorial for continuous-response mixture
  models: Frischkorn and Popov
  (\citeproc{ref-frischkorn2025tutorial}{2025})
\end{itemize}

\section*{References}\label{references}
\addcontentsline{toc}{section}{References}

\phantomsection\label{refs}
\begin{CSLReferences}{1}{0}
\bibitem[\citeproctext]{ref-batchelder1999theoretical}
Batchelder, W. H., \& Riefer, D. M. (1999). Theoretical and empirical
review of multinomial process tree modeling. \emph{Psychonomic Bulletin
\& Review}, \emph{6}(1), 57--86.
\url{https://doi.org/10.3758/BF03210812}

\bibitem[\citeproctext]{ref-burkner2017brms}
Bürkner, P.-C. (2017). {brms}: {A}n {R} package for {B}ayesian
multilevel models using {S}tan. \emph{Journal of Statistical Software},
\emph{80}(1), 1--28. \url{https://doi.org/10.18637/jss.v080.i01}

\bibitem[\citeproctext]{ref-carpenter2017stan}
Carpenter, B., Gelman, A., Hoffman, M. D., Lee, D., Goodrich, B.,
Betancourt, M., Brauer, M., Guo, J., Li, P., \& Riddell, A. (2017).
Stan: {A} probabilistic programming language. \emph{Journal of
Statistical Software}, \emph{76}(1), 1--32.
\url{https://doi.org/10.18637/jss.v076.i01}

\bibitem[\citeproctext]{ref-chalmers2020simdesign}
Chalmers, R. P., \& Adkins, M. C. (2020). Writing effective and reliable
{M}onte {C}arlo simulations with the {SimDesign} package. \emph{The
Quantitative Methods for Psychology}, \emph{16}(4), 248--280.
\url{https://doi.org/10.20982/tqmp.16.4.p248}

\bibitem[\citeproctext]{ref-erdfelder2009multinomial}
Erdfelder, E., Auer, T.-S., Hilbig, B. E., Aßfalg, A., Moshagen, M., \&
Nadarevic, L. (2009). Multinomial processing tree models: {A} review of
the literature. \emph{Zeitschrift f{ü}r Psychologie / Journal of
Psychology}, \emph{217}(3), 108--124.
\url{https://doi.org/10.1027/0044-3409.217.3.108}

\bibitem[\citeproctext]{ref-frischkorn2025tutorial}
Frischkorn, G. T., \& Popov, V. (2025). A tutorial for estimating
{B}ayesian hierarchical mixture models for visual working memory tasks:
{I}ntroducing the {B}ayesian {M}easurement {M}odeling (bmm) package for
{R}. \emph{Behavior Research Methods}, \emph{57}(5), 144.
\url{https://doi.org/10.3758/s13428-025-02643-0}

\bibitem[\citeproctext]{ref-gabry2019visualization}
Gabry, J., Simpson, D., Vehtari, A., Betancourt, M., \& Gelman, A.
(2019). Visualization in {B}ayesian workflow. \emph{Journal of the Royal
Statistical Society: Series A (Statistics in Society)}, \emph{182}(2),
389--402. \url{https://doi.org/10.1111/rssa.12378}

\bibitem[\citeproctext]{ref-gelman1992inference}
Gelman, A., \& Rubin, D. B. (1992). Inference from iterative simulation
using multiple sequences. \emph{Statistical Science}, \emph{7}(4),
457--472. \url{https://doi.org/10.1214/ss/1177011136}

\bibitem[\citeproctext]{ref-goettmann2025recovery}
Göttmann, J., Frischkorn, G. T., Oberauer, K., Schaefer, S. B., \&
Schubert, A.-L. (2025). \emph{Modeling individual differences in working
memory: {S}ubject-level parameter recovery within the {M}emory
{M}easurement {M}odel framework ({M3})}.
\url{https://doi.org/10.31234/osf.io/945d2}

\bibitem[\citeproctext]{ref-gronau2017bridge}
Gronau, Q. F., Sarafoglou, A., Matzke, D., Ly, A., Boehm, U., Marsman,
M., Leslie, D. S., Forster, J. J., Wagenmakers, E.-J., \& Steingroever,
H. (2017). A tutorial on bridge sampling. \emph{Journal of Mathematical
Psychology}, \emph{81}, 80--97.
\url{https://doi.org/10.1016/j.jmp.2017.09.005}

\bibitem[\citeproctext]{ref-heck2018treebugs}
Heck, D. W., Arnold, N. R., \& Arnold, D. (2018). {TreeBUGS}: {A}n {R}
package for hierarchical multinomial-processing-tree modeling.
\emph{Behavior Research Methods}, \emph{50}(1), 264--284.
\url{https://doi.org/10.3758/s13428-017-0869-7}

\bibitem[\citeproctext]{ref-jacoby1991process}
Jacoby, L. L. (1991). A process dissociation framework: Separating
automatic from intentional uses of memory. \emph{Journal of Memory and
Language}, \emph{30}(5), 513--541.
\url{https://doi.org/10.1016/0749-596X(91)90025-F}

\bibitem[\citeproctext]{ref-li2025updating}
Li, C., Frischkorn, G. T., \& Oberauer, K. (2025). Updating of
information in working memory: {T}ime course and consequences.
\emph{Cognitive Psychology}, \emph{156}, 101702.
\url{https://doi.org/10.1016/j.cogpsych.2024.101702}

\bibitem[\citeproctext]{ref-li2026filtering}
Li, C., Frischkorn, G. T., \& Oberauer, K. (in press). Can we process
information without encoding it into working memory? \emph{Journal of
Experimental Psychology: Learning, Memory, and Cognition}.

\bibitem[\citeproctext]{ref-loaiza2024active}
Loaiza, V. M., \& Souza, A. S. (2024). Active maintenance in working
memory reinforces bindings for future retrieval from episodic long-term
memory. \emph{Memory \& Cognition}, \emph{52}(8), 1999--2021.
\url{https://doi.org/10.3758/s13421-024-01596-7}

\bibitem[\citeproctext]{ref-luce1959individual}
Luce, R. D. (1959). \emph{Individual choice behavior: {A} theoretical
analysis}. John Wiley \& Sons.

\bibitem[\citeproctext]{ref-mcelreath2020statistical}
McElreath, R. (2020). \emph{Statistical rethinking: {A} {B}ayesian
course with examples in {R} and {S}tan} (2nd ed.). Chapman; Hall/CRC.
\url{https://doi.org/10.1201/9780429029608}

\bibitem[\citeproctext]{ref-oberauer2019bindings}
Oberauer, K. (2019). Working memory capacity limits memory for bindings.
\emph{Journal of Cognition}, \emph{2}(1), 40.
\url{https://doi.org/10.5334/joc.86}

\bibitem[\citeproctext]{ref-oberauer2023measurement}
Oberauer, K. (2023). Measurement models for visual working memory---{A}
factorial model comparison. \emph{Psychological Review}, \emph{130}(3),
841--852. \url{https://doi.org/10.1037/rev0000328}

\bibitem[\citeproctext]{ref-oberauer2019simple}
Oberauer, K., \& Lewandowsky, S. (2019). Simple measurement models for
complex working-memory tasks. \emph{Psychological Review},
\emph{126}(6), 880--932. \url{https://doi.org/10.1037/rev0000159}

\bibitem[\citeproctext]{ref-oberauer2018benchmarks}
Oberauer, K., Lewandowsky, S., Awh, E., Brown, G. D. A., Conway, A.,
Cowan, N., Donkin, C., Farrell, S., Hitch, G. J., Hurlstone, M. J., Ma,
W. J., Morey, C. C., Nee, D. E., Schweppe, J., Vergauwe, E., \& Ward, G.
(2018). Benchmarks for models of short-term and working memory.
\emph{Psychological Bulletin}, \emph{144}(9), 885--958.
\url{https://doi.org/10.1037/bul0000153}

\bibitem[\citeproctext]{ref-bmm2025}
Popov, V., Frischkorn, G. T., Li, C., \& Bürkner, P.-C. (2025).
\emph{Bmm: {E}asy and accessible {B}ayesian measurement models using
'brms'}. \url{https://doi.org/10.32614/CRAN.package.bmm}

\bibitem[\citeproctext]{ref-schad2021toward}
Schad, D. J., Betancourt, M., \& Vasishth, S. (2021). Toward a
principled {B}ayesian workflow in cognitive science. \emph{Psychological
Methods}, \emph{26}(1), 103--126.
\url{https://doi.org/10.1037/met0000275}

\bibitem[\citeproctext]{ref-schubert2023working}
Schubert, A.-L., Löffler, C., Sadus, K., Göttmann, J., Hein, J.,
Schröer, P., Teuber, A., \& Hagemann, D. (2023). Working memory load
affects intelligence test performance by reducing the strength of
relational item bindings and impairing the filtering of irrelevant
information. \emph{Cognition}, \emph{236}, 105438.
\url{https://doi.org/10.1016/j.cognition.2023.105438}

\bibitem[\citeproctext]{ref-talts2018validating}
Talts, S., Betancourt, M., Simpson, D., Vehtari, A., \& Gelman, A.
(2018). \emph{Validating {B}ayesian inference algorithms with
simulation-based calibration}. \url{https://arxiv.org/abs/1804.06788}

\bibitem[\citeproctext]{ref-thurstone1927law}
Thurstone, L. L. (1927). A law of comparative judgment.
\emph{Psychological Review}, \emph{34}(4), 273--286.
\url{https://doi.org/10.1037/h0070288}

\bibitem[\citeproctext]{ref-wilson2019ten}
Wilson, R. C., \& Collins, A. G. E. (2019). Ten simple rules for the
computational modeling of behavioral data. \emph{eLife}, \emph{8},
e49547. \url{https://doi.org/10.7554/eLife.49547}

\bibitem[\citeproctext]{ref-yellott1977relationship}
Yellott, Jr., John I. (1977). The relationship between {L}uce's choice
axiom, {T}hurstone's theory of comparative judgment, and the double
exponential distribution. \emph{Journal of Mathematical Psychology},
\emph{15}(2), 109--144.
\url{https://doi.org/10.1016/0022-2496(77)90026-8}

\bibitem[\citeproctext]{ref-yonelinas2012process}
Yonelinas, A. P., \& Jacoby, L. L. (2012). The process-dissociation
approach two decades later: Convergence, boundary conditions, and new
directions. \emph{Memory \& Cognition}, \emph{40}(5), 663--680.
\url{https://doi.org/10.3758/s13421-012-0205-5}

\end{CSLReferences}






\end{document}
